% =============================================================================
% Protected twisted-M obstruction and admissible local reduction.
%
% This fragment isolates the exact obstruction datum whose vanishing
% promotes Costello's twisted-M-theory motivation to a protected
% Hamiltonian realization at a brane point.
% =============================================================================

\subsection{Protected twisted-M motivation and local reduction}\label{ssec:twistedM-universality}

Perturbative twisted M-theory in the sense of
Costello~\cite{costello-twistedM} enters through a protected
Hamiltonian realization datum at a brane point.  The mathematical
assertion is a realization criterion: vanishing of the period,
closed-field comparison, scalar-anomaly, and brane-defect QME
obstructions constructs the admissible Hamiltonian sector to which the
formal local theorem of \S\ref{ssec:local-theorem} applies.

\begin{defn}[Protected Hamiltonian sector at a brane point]
  Let $(X,\omega)$ be a connected holomorphic-symplectic surface, let
  $p\in X$, and let
  $\mc T_{\mathrm{tw}M}\bigl(\R\times X;\R\times p\bigr)$ denote
  perturbative twisted-M-theory data on $\R\times X$ in the presence
  of a stack of $N$ one-dimensional topological branes supported on
  $\R\times p$.  A \emph{protected Hamiltonian sector at $\widehat p$}
  is a restricted datum
  $$
  \mc T_{\mathrm{tw}M}\bigl(\R\times X;\R\times p\bigr)
    \big|_{\mathrm{Ham},\widehat p}
    =
    \bigl(
      \mc E_{\mathrm{open},\mathrm{prot}},
      \mc E_{\mathrm{cl},\mathrm{prot}},
      \partial_{\mathrm{bb}}^{\mathrm{prot}}
    \bigr)
  $$
  consisting of the protected open BV theory on the brane line, the
  protected closed BV theory coming from formal Hamiltonian operators
  at the symplectic disk $\widehat X_p$, and the protected
  bulk-to-boundary evaluation on the brane algebra.

  The realization problem is to construct this triple from the ambient
  BV fields and to prove the vanishing of the obstruction datum below.
\end{defn}

\begin{defn}[Admissible protected Hamiltonian restriction]\label{defn:admissible-protected-hamiltonian-restriction}
  A protected Hamiltonian sector is \emph{admissible for the local
  theorem} if it comes with the following data.
  \begin{enumerate}[label=(\roman*)]
    \item \emph{Formal Darboux Hamiltonian reduction.}  The completed
      closed Hamiltonian algebra at $p$ is identified, up to formal
      symplectomorphism, with
      $\widehat{\mathfrak h}=\C[[z_1,z_2]]/\C\cdot 1$, and the point
      brane Ext algebra is identified with
      $\Ext_X^*(\mc O_p,\mc O_p)\cong\Lambda^*T_pX$, with cyclic trace
      induced by $\omega_p$.
    \item \emph{Brane-defect locality.}  The protected operator
      restriction to $\R\times p$ factors through the weighted Tate
      field complex and is compatible with the interval-wise
      factorization map
      $\Phi^{\mathrm{Ham}}_{\mathrm{fact}}$, the chain-level
      primitive projection, and the weighted Hadamard parametrix used
      in the local theorem.
    \item \emph{Scalar anomaly compatibility.}  The centre-of-mass
      scalar anomaly of the protected brane stack maps to the class
      $[\bar c]\in H^2_{\mathrm{Lie}}(\bar A,\C)$ used in the local
      scalar-anomaly comparison.
  \end{enumerate}
\end{defn}

\begin{defn}[Protected twist obstruction datum]\label{defn:protected-twist-obstruction-datum}
  Fix a candidate protected twisted-M background near
  \(\R\times p\) and a candidate comparison with the Hamiltonian
  formal-disk model.  The comparison carries four obstruction classes.

  First, every holomorphic symplectic vector field \(v\) determines the
  closed one-form
  \[
    \alpha_v=-\iota_v\omega.
  \]
  The holomorphic de Rham exact sequence
  \[
    0\longrightarrow\C_X\longrightarrow\mc O_X
      \xrightarrow{d}\Omega^1_{X,\mathrm{cl}}\longrightarrow0
  \]
  gives the period obstruction
  \[
    \operatorname{per}(v)=\delta(\alpha_v)\in H^1(X,\C_X).
  \]
  This is the obstruction to choosing a single global holomorphic
  Hamiltonian primitive for \(v\).

  Second, a chosen local comparison of closed protected fields with the
  Hamiltonian BF fields gives a cone
  \[
  \mc C_{\mathrm{cl},p}=
  \operatorname{Cone}\!\left(
    \mc E_{\mathrm{cl},\mathrm{prot},c}^{!}(\R\times\widehat X_p)
      \xrightarrow{\ \chi_{\mathrm{cl}}\ }
    \Omega^\bullet_c(\R)\widehat\otimes
      (\widehat{\mathfrak h}\ltimes
       \widehat{\mathfrak h}^{\vee}_{\mathrm{cont}}[1])^{!}
  \right).
  \]
  The comparison is a field-level equivalence only after this cone is
  acyclic in the protected summand.

  Third, the scalar anomaly mismatch is
  \[
    \operatorname{ob}_{\mathrm{scal}}
      =
    \chi_{\mathrm{anom}}\bigl([\operatorname{anom}_{\mathrm{prot}}]\bigr)
      -[\bar c]\in H^2_{\mathrm{Lie}}(\bar A,\C),
  \]
  where \(\chi_{\mathrm{anom}}\) is the anomaly map induced by the
  comparison and \([\bar c]\) is the scalar class used in the local
  theorem.

  Fourth, after choosing the weighted coefficient model used by the
  local theorem, the brane-defect quantum obstruction is the class
  \[
  \operatorname{ob}_{\mathrm{QME}}=
  \left[
    QI_{\mathrm{prot}}
    +\frac12\{I_{\mathrm{prot}},I_{\mathrm{prot}}\}
    +\hbar\Delta I_{\mathrm{prot}}
  \right]
  \in
  H^1\!\left(
    \mathcal O_{\mathrm{loc}}(\mc E_{\mathrm{prot},w}),
    Q+\{I_{0,\mathrm{prot}},-\}
  \right).
  \]
  Its vanishing is the precise statement that the protected brane-defect
  interaction satisfies the QME in the same normalization as
  Problem~\ref{prob:weighted-rg-locality}.
\end{defn}

\begin{defn}[Super-balanced equivariant QME datum]\label{defn:super-balanced-equivariant-qme-datum}
  A super-balanced equivariant brane-defect datum consists of:
  \begin{enumerate}[label=(\roman*)]
    \item a matrix Lie superalgebra \(\mathfrak s\) acting on the
      Chan-Paton factor, together with an even nondegenerate invariant
      form \(B_{\mathfrak s}\) used to build the BV heat-kernel
      contraction;
    \item a trace character \(\operatorname{tr}_{\mathfrak s}\) on the
      matrix factor and a lifted identity element \(I_{\mathfrak s}\)
      with
      \[
        \operatorname{tr}_{\mathfrak s}(I_{\mathfrak s})=0;
      \]
    \item the weighted brane-defect local-functional complex
      \[
        \mathfrak Q_{\mathfrak s}=
        \left(
        \mathcal O_{\mathrm{loc}}(\mc E^{\mathfrak s}_{\mathrm{bd},w}),
        d_{\mathrm{QME}}=Q+\{I_{0,\mathfrak s},-\}
        \right);
      \]
    \item a \(P_{(z_1)}\)-equivariant coefficient topology on the formal
      symplectic disk, where
      \[
        P_{(z_1)}=\{\varphi\in
        \operatorname{Symp}_{\mathrm{form}}(\widehat{\C^2}_0):
        \varphi^*(z_1)\in (z_1)\}.
      \]
  \end{enumerate}
  The scalar associated graded of \(\mathfrak Q_{\mathfrak s}\) is the
  Chevalley--Eilenberg complex of the reduced Hamiltonian algebra
  \(\bar A=\C[z_1,z_2]/\C\cdot1\) with coefficients in the central
  boundary ghost line \(\C\gamma_{\mathbf 1}\).
\end{defn}

\begin{thm}[Theorem F double-prime: super-balanced equivariant scalar QME vanishing]\label{thm:joint-super-balanced-symp-qme-vanishing}
  Let
  \((\mathfrak s,B_{\mathfrak s},\operatorname{tr}_{\mathfrak s})\)
  be a super-balanced equivariant datum in the sense of
  Definition~\ref{defn:super-balanced-equivariant-qme-datum}.  For
  \(a,b\in\Omega^0_c(\R)\) and
  \(f,g\in\C[z_1,z_2]\), the scalar part of the mixed brane-defect QME
  obstruction is represented in \(\mathfrak Q_{\mathfrak s}\) by
  \[
    \operatorname{Ob}^{\mathfrak s}_{\mathrm{sc}}
    (\gamma_{\mathbf 1};a,f;b,g)
    =
    \hbar\,\operatorname{tr}_{\mathfrak s}(I_{\mathfrak s})\,
    \omega(f,g)
    \int_{\R}a(t)b(t)\gamma_{\mathbf 1}(t)\,dt .
  \]
  Hence the scalar obstruction vanishes identically for a
  super-balanced datum.  The vanishing is \(P_{(z_1)}\)-equivariant.
  Equivalently, the associated graded scalar class is
  \[
    [\operatorname{Ob}^{\mathfrak s}_{\mathrm{sc}}]
    =
    \hbar\,\operatorname{tr}_{\mathfrak s}(I_{\mathfrak s})\,
    [\bar c]
    \in
    H^2_{\mathrm{Lie}}(\bar A,\C\gamma_{\mathbf 1}),
  \]
  and the theorem applies exactly on the super-balanced locus
  \(\operatorname{tr}_{\mathfrak s}(I_{\mathfrak s})=0\).
  Thus the scalar channel is an accepted universal Koszul-duality input
  precisely on this locus: its acceptance component is the class
  \(\hbar\,\operatorname{tr}_{\mathfrak s}(I_{\mathfrak s})[\bar c]\).

  If, in addition, the reduced non-scalar summand
  \[
    \mathfrak Q_{\mathfrak s}^{\mathrm{red}}
      =
      \operatorname{Cone}\!\left(
      \mathfrak Q_{\mathfrak s}
      \longrightarrow
      C^\bullet_{\mathrm{Lie}}(\bar A;\C\gamma_{\mathbf 1})
      \right)
  \]
  has zero degree-one cohomology after the chosen weighted regulator and
  interval-factorization null-homotopies are imposed, then the complete
  degree-one brane-defect QME obstruction class in
  \(H^1(\mathfrak Q_{\mathfrak s})\) vanishes.
\end{thm}

\begin{proof}
  The scalar contraction in the BV bracket is obtained by closing the
  matrix index in the single boundary loop.  For an ordinary matrix
  trace this contraction is \(N\).  For a supertrace datum it is the
  trace of the lifted identity, namely
  \(\operatorname{tr}_{\mathfrak s}(I_{\mathfrak s})\).  The target
  part of the same local bracket is the constant Taylor coefficient
  cocycle
  \[
    \omega(f,g)=[\{f,g\}]_0
  \]
  of Lemma~\ref{lem:omega-cocycle}, and the brane-line locality gives
  the compactly supported factor
  \(\int_\R a(t)b(t)\gamma_{\mathbf 1}(t)\,dt\).  This proves the
  displayed representative.

  Passing to the scalar associated graded keeps only the Hamiltonian
  cocycle \([\bar c]\) and the matrix coefficient
  \(\operatorname{tr}_{\mathfrak s}(I_{\mathfrak s})\), giving the
  displayed cohomology class.  Super-balance sets this coefficient to
  zero, so the degree-one representative is the zero cochain, not
  merely an exact cochain.  The \(P_{(z_1)}\)-equivariance follows because
  \(\operatorname{tr}_{\mathfrak s}(I_{\mathfrak s})\) is matrix-factor
  data and \(\omega\) is the symplectic Poisson cocycle on the target
  factor.  The completion at the brane axis is functorial only under
  \(P_{(z_1)}\), which is exactly the equivariance group in the
  hypothesis.

  The last assertion is the long exact sequence of cohomology for the
  displayed cone.  The scalar quotient class is zero by the first part,
  and the reduced degree-one cohomology is zero by hypothesis.  Therefore
  the total degree-one obstruction class in
  \(H^1(\mathfrak Q_{\mathfrak s})\) vanishes.
\end{proof}

\begin{defn}[Spin-Hecke queer coefficient input]\label{defn:spin-hecke-queer-input}
  Let \(N\ge1\) and
  \[
    \mathfrak q(N)=
    \left\{
      \begin{pmatrix} A&B\\ B&A\end{pmatrix}:
      A,B\in\mathfrak{gl}_N
    \right\}\subset\mathfrak{gl}(N|N).
  \]
  Put
  \[
    J=\begin{pmatrix}0&I_N\\ -I_N&0\end{pmatrix},
    \qquad
    P^{\mathfrak q}=\begin{pmatrix}I_N&0\\0&-I_N\end{pmatrix},
  \]
  and define the queer trace
  \[
    \operatorname{otr}
      \begin{pmatrix}A&B\\ B&A\end{pmatrix}
      =\operatorname{Tr}(B)\in\Pi\C .
  \]
  A spin-Hecke input for the queer brane is the Sergeev
  Hecke--Clifford action on \(V^{\otimes n}\), \(V=\C^{N|N}\), whose
  central-character comparison identifies
  \[
    \operatorname{Tr}_{\mathfrak{gl}_N}(I_N)
      =
    \operatorname{otr}_{\mathfrak q(N)}(J)
      =N .
  \]
\end{defn}

\begin{thm}[Theorem G odd-trace: queer centre anomaly from the spin-Hecke input]\label{thm:queer-u1-center-anomaly-spin-hecke}
  Assume the spin-Hecke input of
  Definition~\ref{defn:spin-hecke-queer-input}.  For the queer-trace
  boundary evaluation
  \(f\mapsto\operatorname{otr} f(\phi_1,\phi_2)\), the odd scalar QME
  representative is
  \[
    \operatorname{Ob}^{\mathrm{otr}}_{\mathrm{sc}}
    (\gamma_{\mathbf 1};a,f;b,g)
    =
    \hbar N\,\omega(f,g)
    \int_\R a(t)b(t)\gamma_{\mathbf 1}(t)\,dt
    \in
    \mathcal O_{\mathrm{loc}}(\mc E^{\mathfrak q}_{\mathrm{bd},w})_{\bar1}.
  \]
  Its associated graded CE class is
  \[
    \hbar N[\bar c]^{\mathrm{otr}}
    \in H^2_{\mathrm{Lie}}(\bar A,\Pi\C)_{\bar1}.
  \]
  This class is nonzero and is independent of the even class
  \(\hbar N[\bar c]\in H^2_{\mathrm{Lie}}(\bar A,\C)_{\bar0}\).
  Since \(N\ge1\), the class does not vanish over \(\C[[\hbar]]\) before
  imposing an odd scalar quotient or a target null-homotopy of the odd
  line \(\Pi\C\cdot[\bar c]^{\mathrm{otr}}\).
  Hence the odd-trace channel is accepted by a universal Koszul-duality
  target exactly after that quotient or null-homotopy is included in the
  target datum; otherwise its acceptance obstruction is the displayed
  nonzero odd class.
  The spin-Hecke input justifies the coefficient identity at the
  cohomology layer.  It does not give a chain-level QME cancellation:
  \[
    P^{\mathfrak q}J(P^{\mathfrak q})^{-1}=-J
  \]
  is the obstruction to the unsigned parity-equivariant regulator used
  in the super-balanced theorem above.
\end{thm}

\begin{proof}
  The queer brane has two central coefficient lines:
  the even line generated by the identity and the odd line generated by
  \(J\).  The queer trace kills the even identity and evaluates the odd
  central element by
  \[
    \operatorname{otr}(J)=
    \operatorname{Tr}(I_N)=N .
  \]
  The spin-Hecke input asserts that this coefficient is preserved by the
  Sergeev central-character comparison; the Hamiltonian cocycle on the
  target remains the same \(\omega(f,g)=[\{f,g\}]_0\).  Therefore the
  scalar boundary loop gives exactly the displayed odd representative.

  Passing to associated graded removes the brane-line smearing and keeps
  the Lie two-cocycle \(\Pi\omega\), hence the class
  \(\hbar N[\bar c]^{\mathrm{otr}}\).  Nontriviality is the same
  calculation as Lemma~\ref{lem:omega-cocycle}: a primitive
  \(\bar\eta:\bar A\to\Pi\C\) would lift with
  \(\tilde\eta(1)=0\), but then
  \[
    \delta\tilde\eta(z_1,z_2)
      =-\tilde\eta(\{z_1,z_2\})
      =-\tilde\eta(1)=0
  \]
  while \(\omega(z_1,z_2)=1\).
  The even and odd classes are independent because they lie in the two
  direct summands
  \(H^2_{\mathrm{Lie}}(\bar A,\C)\oplus
  H^2_{\mathrm{Lie}}(\bar A,\Pi\C)\).
  Since \(N\ge1\) in the datum and \(\hbar\) is formal, the displayed
  odd class is not killed inside
  \(H^2_{\mathrm{Lie}}(\bar A,\Pi\C)[[\hbar]]\).  Killing it requires an
  additional target quotient or null-homotopy; it is not a consequence
  of the spin-Hecke coefficient comparison.

  Finally,
  \(P^{\mathfrak q}J(P^{\mathfrak q})^{-1}=-J\) is an explicit matrix
  computation.  Thus the queer channel has signed parity, not the
  unsigned parity condition used in the super-balanced cancellation.
  The algebraic spin-Hecke input gives the coefficient class, but not a
  null-homotopy of this odd obstruction.
\end{proof}

\begin{lem}[Primary obstruction test]\label{lem:protected-primary-obstruction-test}
  A protected twisted-M background can supply an admissible protected
  Hamiltonian restriction only if the period class, the comparison cone,
  the scalar mismatch, and the QME obstruction in
  Definition~\ref{defn:protected-twist-obstruction-datum} vanish in the
  senses stated there.  Conversely, if the period classes vanish for the
  Hamiltonian vector fields under consideration, the comparison cone is
  acyclic on the protected summand, the scalar mismatch is zero, and a
  solution of the displayed QME class has been chosen compatibly with
  interval factorization, then the candidate comparison supplies the data
  of Definition~\ref{defn:admissible-protected-hamiltonian-restriction}.
\end{lem}

\begin{proof}
  The period statement is the connecting morphism of
  \(0\to\C_X\to\mc O_X\xrightarrow d\Omega^1_{X,\mathrm{cl}}\to0\):
  local Hamiltonian primitives \(h_i\) for \(\alpha_v\) differ by
  constants on overlaps, and that Cech cocycle is exactly
  \(\delta(\alpha_v)\).  It vanishes precisely when the local primitives
  glue to a global Hamiltonian, modulo one global constant.

  A comparison of protected closed fields with Hamiltonian BF fields is a
  quasi-isomorphism on the protected summand exactly when its mapping cone
  is acyclic.  If the cone has cohomology, the unmatched class is a
  protected closed field invisible to the formal Hamiltonian model, or a
  Hamiltonian BF field not produced by the protected sector.

  Scalar compatibility is an equality in
  \(H^2_{\mathrm{Lie}}(\bar A,\C)\), since both the protected anomaly and
  the local scalar anomaly are central Lie algebra cocycles after the
  comparison.  Finally, the displayed QME expression is the standard
  degree-one obstruction class for deforming the classical interaction to
  a quantum BV interaction in the fixed weighted coefficient category.
  The interval-factorization compatibility is precisely the
  brane-defect locality clause in the admissibility definition.
\end{proof}

\begin{defn}[Ambient-to-protected realization map]\label{defn:ambient-protected-realization-map}
  An ambient-to-protected realization map is a morphism from the ambient
  twisted-M BV fields near \(\R\times p\) to the triple
  \[
    \bigl(
      \mc E_{\mathrm{open},\mathrm{prot}},
      \mc E_{\mathrm{cl},\mathrm{prot}},
      \partial_{\mathrm{bb}}^{\mathrm{prot}}
    \bigr)
  \]
  whose closed-field cone is \(\mc C_{\mathrm{cl},p}\), whose scalar
  anomaly map is \(\chi_{\mathrm{anom}}\), and whose brane interaction
  has QME obstruction \(\operatorname{ob}_{\mathrm{QME}}\).
\end{defn}

\begin{prop}[Admissible protected-sector reduction]
  Let $(X,\omega)$ be a connected holomorphic-symplectic surface and
  let $p\in X$.  Suppose that a perturbative twisted-M background with
  branes on $\R\times p$ supplies an admissible protected Hamiltonian
  sector at $\widehat p$.  Then, after the Darboux coordinate and
  scalar-anomaly normalizations included in the admissibility data, that
  restricted sector is identified with the formal local open--closed model
  of \S\ref{ssec:local-theorem}.

  The open side is
  $$
  \Omega^*(\R)\otimes
  \Lambda(\varepsilon_1,\varepsilon_2)\otimes\lie{gl}_N[1],
  $$
  the closed side is
  $$
  C^\bullet_{\mathrm{CE,cont}}(\mathfrak g),
  \qquad
  \mathfrak g=
  \widehat{\mathfrak h}\ltimes
  \widehat{\mathfrak h}^{\vee}_{\mathrm{cont}}[1],
  \qquad
  \widehat{\mathfrak h}=\C[[z_1,z_2]]/\C\cdot 1,
  $$
  and the bulk-to-boundary evaluation is
  $$
  f\longmapsto
  \overline{\mathrm{Tr}\,f(\phi_1,\phi_2)}.
  $$
  Hence the proved Hamiltonian, factorization, and scalar-anomaly
  statements of the local theorem apply to the protected sector.  The
  all-order one-antifield quantum extension is still conditional on
  vanishing of the mixed brane-defect QME obstruction of
  Problem~\ref{prob:weighted-rg-locality} for that sector.
\end{prop}

\begin{proof}
  The formal Darboux part of admissibility identifies the completed
  Hamiltonian algebra at the brane point with
  $\C[[z_1,z_2]]/\C\cdot 1$, after choosing formal Darboux
  coordinates.  The point-brane Ext algebra and its cyclic trace give
  the open coefficient data of the Dirac probe model.

  Brane-defect locality supplies the factorization-algebraic
  restriction to the brane line and matches it with the local
  factorization map, primitive projection, and weighted Hadamard
  reduction already used in the local theorem.  These
  hypotheses put the protected sector in the same algebraic category
  as the local model.  QME-flatness for the one-antifield quantum
  extension is the vanishing condition
  \(\operatorname{ob}_{\mathrm{QME}}=0\) in
  Theorem~\ref{thm:protected-twist-hamiltonian-criterion}.

  Scalar anomaly compatibility identifies the brane centre-of-mass
  anomaly with the class $[\bar c]$ appearing in the local anomaly
  comparison.  The three admissibility inputs therefore supply exactly
  the data on which the local theorem acts.  Applying that theorem
  gives the asserted protected-sector reduction.
\end{proof}

\begin{thm}[Protected twist-to-Hamiltonian comparison criterion]\label{thm:protected-twist-hamiltonian-criterion}
  Let \((X,\omega)\) be a connected holomorphic-symplectic surface and
  \(p\in X\).  Suppose a perturbative twisted-M background with branes on
  \(\R\times p\) supplies a candidate protected sector together with the
  obstruction datum of
  Definition~\ref{defn:protected-twist-obstruction-datum}.  Assume:
  \begin{enumerate}[label=(\roman*)]
    \item the period classes \(\operatorname{per}(v)\) vanish for the
      locally Hamiltonian symplectic vector fields used by the sector;
    \item the closed-field comparison cone \(\mc C_{\mathrm{cl},p}\) is
      acyclic on the protected summand;
    \item \(\operatorname{ob}_{\mathrm{scal}}=0\);
    \item \(\operatorname{ob}_{\mathrm{QME}}=0\), with a chosen solution
      compatible with compactly supported interval factorization.
  \end{enumerate}
  Then the candidate protected sector determines an admissible protected
  Hamiltonian restriction.  Consequently the formal local Hamiltonian
  theorem applies to that sector.  If, in addition, the QME solution
  includes the one-antifield weighted Tate lift of
  Problem~\ref{prob:weighted-rg-locality}, then the corresponding
  one-antifield quantum Hamiltonian comparison is valid for that
  protected sector and for that sector only.
\end{thm}

\begin{proof}
  By Lemma~\ref{lem:protected-primary-obstruction-test}, the four
  hypotheses supply the Darboux Hamiltonian reduction, the field-level
  protected comparison, scalar anomaly compatibility, and brane-defect
  QME locality required in
  Definition~\ref{defn:admissible-protected-hamiltonian-restriction}.
  Applying Proposition~\ref{prop:formal-local-global-restriction} at the
  formal disk and then the admissible protected-sector reduction gives the
  local Hamiltonian BF model and its bulk-to-boundary evaluation.  The last
  sentence is exactly the conditional weighted Tate/QME extension isolated in
  Problem~\ref{prob:weighted-rg-locality}, transported through the chosen
  protected-sector equivalence.
\end{proof}

\begin{lem}[Formal-restriction external acceptance criterion]\label{lem:admissibility-not-globalization}
  An admissible protected Hamiltonian restriction contributes precisely
  the formal Hamiltonian disk, the Dirac brane Ext algebra, the brane-line
  factorization restriction, and the scalar anomaly class \( [\bar c] \).
  For an external compact BCOV, CY-to-chiral, Igusa, Borcherds, or BKM
  target \(T\), this contribution is an accepted universal Koszul-duality
  input exactly after the target supplies its convention datum
  \(\mathfrak C_T\), a morphism
  \(\mathfrak C_{\widehat p}^{\mathrm{Ham}}\to \mathfrak C_T\), and
  null-homotopies for the obstruction vector
  \[
    \operatorname{Ob}^{\mathrm{ext}}_T
    =
    \bigl(
      \operatorname{per}_T,
      \operatorname{ob}^{\mathrm{desc}}_T,
      \operatorname{ob}^{\mathrm{BCOV}}_T,
      \operatorname{ob}^{\mathrm{VolIII}}_T,
      \operatorname{ob}^{\mathrm{Igusa/BKM}}_T
    \bigr)
    \in H^1(\mathfrak K_T),
  \]
  with a component omitted only when \(T\) does not assert the
  corresponding structure.  A non-zero component is the exact obstruction
  to the corresponding external theorem.
\end{lem}

\begin{proof}
  The three admissibility clauses record Darboux Hamiltonians at
  \(\widehat X_p\), the point-brane Ext algebra, brane-defect locality on
  \(\R\times p\), and scalar anomaly compatibility.  They contain no
  compact Calabi--Yau threefold, no negative-cyclic Calabi--Yau class,
  no Vol III specialization datum
  \(\Phi_d=\mathrm{Sp}_{\Sigma_{d-1},C}\circ\Phi^{\mathrm{FA}}_d\), no
  compact worldsheet or \(S^3\) framing, no BCOV propagator or
  holomorphic anomaly equation, no Borcherds input Jacobi form, no Weyl
  chamber, and no compact Hall orientation.  Two ambient theories can
  have the same admissible formal restriction and different choices of
  all those global data.  In the target convention complex
  \(\mathfrak K_T\) these choices are independent coordinates of
  \(\operatorname{Ob}^{\mathrm{ext}}_T\).  Vanishing with specified
  null-homotopies is therefore precisely the acceptance datum for using
  the formal restriction in \(T\); a non-zero coordinate is precisely the
  named target obstruction.
\end{proof}

\begin{defn}[Protected-sector matched-conventions contribution]\label{defn:protected-sector-matched-conventions-contribution}
  The protected-sector reduction contributes the following datum to any
  external matched-conventions theorem:
  \[
  \mathfrak C_{\widehat p}^{\mathrm{Ham}}=
  \bigl(
    d=2,\widehat X_p,\omega_p,
    \R_{\mathrm{brane}},
    \widehat{\mathfrak h}=\C[[z_1,z_2]]/\C,
    \Ext_X^*(\mc O_p,\mc O_p),
    [\bar c],
    \Phi^{\mathrm{Ham}}_{\mathrm{fact}}
  \bigr).
  \]
  The remaining datum for a compact BCOV, CY-to-chiral, or Igusa/BKM
  theorem consists of the compact target, its Calabi--Yau or modular
  class, the worldsheet or specialization curve, the framing, the global
  propagator or Hall orientation, and the target obstruction complex.
\end{defn}

\begin{cor}[Formal uniformity after admissibility]\label{cor:protected-formal-uniformity}
  For every connected holomorphic-symplectic surface \((X,\omega)\), every
  brane point \(p\in X\), and every protected twisted-M realization whose
  obstruction datum vanishes, the restricted sector is identified with the
  same formal Hamiltonian brane model
  \[
    \bigl(\widehat{\mathfrak h}\ltimes
    \widehat{\mathfrak h}^{\vee}_{\mathrm{cont}}[1],
    \Omega^*(\R)\otimes\Lambda(\varepsilon_1,\varepsilon_2)
      \otimes\lie{gl}_N[1],
    [\bar c]\bigr)
  \]
  after choosing a formal Darboux coordinate at \(p\).
\end{cor}

\begin{proof}
  The period vanishing gives global Hamiltonian primitives for the vector
  fields used in the sector.  The acyclic closed comparison cone identifies
  the protected closed fields with
  \(\widehat{\mathfrak h}\ltimes
  \widehat{\mathfrak h}^{\vee}_{\mathrm{cont}}[1]\).  The Ext-algebra
  identification gives the Dirac brane probe, and scalar compatibility
  fixes the same class \([\bar c]\).  These are exactly the inputs of
  Proposition~\ref{prop:formal-local-global-restriction}.
\end{proof}

\begin{rmk}[Dirac language inside the protected sector]
  For an admissible protected Hamiltonian sector, the Dirac-thesis
  language of Remark~\ref{rmk:dirac-thesis} is a reading of the
  already identified formal local model.  The brane stack probes the
  formal symplectic normal disk; Dirac reduction gives the open trace
  algebra; the closed Hamiltonian BF theory is the shifted cotangent
  theory of canonical transformations of that disk; and the closed
  observables identify with derived Poisson-central operations on the
  stable open Hamiltonian brane algebra only in the reduced local sense
  proved by Theorem~\ref{thm:main-local}.  The scalar anomaly is the
  same class $[\bar c]$ on the closed and open sides.  The corresponding
  ambient centrality statement is the matched-conventions interface of
  Corollary~\ref{cor:protected-matched-conventions-interface}.
\end{rmk}

\begin{cor}[Protected-sector matched-conventions acceptance criterion]\label{cor:protected-matched-conventions-interface}
  A theorem in a compact BCOV, CY-to-chiral, or Igusa/BKM target may use
  the protected-sector reduction exactly by a morphism of convention data
  from \(\mathfrak C_{\widehat p}^{\mathrm{Ham}}\) to its target datum,
  together with a proof that the target period, descent, anomaly, and
  operadic obstruction classes vanish.
\end{cor}

\begin{proof}
  Lemma~\ref{lem:admissibility-not-globalization} identifies the
  external acceptance class
  \(\operatorname{Ob}^{\mathrm{ext}}_T\in H^1(\mathfrak K_T)\).  A
  morphism of convention data and null-homotopies for this class supply
  exactly the compact, worldsheet, modular, orientation, and descent
  coordinates needed by the target theorem.  Corollary~\ref{cor:protected-formal-uniformity}
  then supplies the \(d=2\) formal-disk normalization functorially in the
  chosen Darboux chart.
\end{proof}

% --- End of fragment.
